\chapter{Ordinary representations and characters}
\label{chap:library:characters}

Pullback, inflation and tensoring with a linear character relate
representations of different groups or give new representations of the
same group. Their character formulas follow from the representing
operators. We first show how characters determine irreducible
representations, then study these operations and extension across a
cyclic quotient.

Throughout, $K$ is a field of characteristic zero and $G$ is a finite group.
A representation $\rho$ of $G$ on a finite dimensional $K$-vector space $V$
has character $\chi_\rho(g)=\tr_K(\rho(g))$. We state algebraic closure
explicitly where it is used. Character separation for finite groups
also holds over an arbitrary field of characteristic zero, as the
averaging argument below shows. Modular systems and reduction are
introduced in Chapter~\ref{chap:library:reduction}.

\section{Representations and their characters}
\label{lib:ordinary:realisations}

\begin{definition}
A representation of $G$ on $V$ is a homomorphism
$\rho:G\to\GL_K(V)$. It is irreducible if $V\ne0$ and its only invariant
subspaces are $0$ and $V$. An equivalence $\rho\simeq\sigma$ is a linear
isomorphism $T:V\to W$ satisfying
\[
 T\rho(g)=\sigma(g)T\qquad(g\in G).
\]
The set $\Irr_K(G)$ consists of the functions $\chi:G\to K$ afforded by
finite dimensional irreducible $K$-representations.
\end{definition}

Thus equality in $\Irr_K(G)$ is equality at every group element. Choosing
different bases or isomorphic realisations does not create new characters:
conjugate operators have equal traces. The value at the identity is
$\chi_\rho(1)=\dim_K V$, regarded as an element of $K$. Every character is
constant on conjugacy classes, since $\rho(xgx^{-1})$ is conjugate to
$\rho(g)$.

\begin{implementation}
Mathlib supplies \lean{Representation},
\lean{Representation.IsIrreducible} and the trace function
\lean{Representation.character}, together with the identities
\lean{Representation.char_one}, \lean{Representation.char_conj} and
\lean{Representation.char_iso}.
The additional type \lean{ModularRep.OrdinaryIrreducibleCharacter.Irr}
collects the functions admitting an irreducible realisation on some
$K^n$, expressed by
\lean{ModularRep.OrdinaryIrreducibleCharacter.Realisation} using these
Mathlib definitions. Equality of these character functions is pointwise,
as expressed by
\lean{ModularRep.OrdinaryIrreducibleCharacter.ext}.
\end{implementation}

\begin{proposition}[Characters determine irreducible representations]
\label{lib:ordinary:separation}
Assume that $K$ is algebraically closed. If two finite dimensional
irreducible $K$-representations $\rho$ and $\sigma$ have equal characters,
then $\rho\simeq\sigma$.
\end{proposition}
For the corresponding statement about traces of simple modules over a
finite dimensional algebra, see, for example,
\cite[Theorem~1.19]{Navarro1998}.
\begin{proof}
Regard the two representation spaces as simple $KG$-modules. Equality of
their traces on $G$ implies equality on every element of $KG$ by linearity.
Suppose that the simple modules are nonisomorphic. Schur's lemma says that
their mutual homomorphism spaces are zero and their endomorphisms are
scalars. The density theorem applied to their direct sum therefore supplies
an element of $KG$ acting as an operator of trace one on the first module
and as zero on the second. This contradicts equality of traces.

An operator of trace one can be obtained by choosing a basis and taking the
projection onto one basis vector. Using this operator also makes the same
separation argument valid in positive characteristic, where the trace of
the identity may be zero.
\end{proof}

\begin{implementation}
The theorem is
\lean{ModularRep.ModularTraceSeparation.representation_equiv_of_character_eq}.
Mathlib supplies the Schur and density statements used in the proof,
\lean{IsSimpleModule.algebraMap_end_bijective_of_isAlgClosed} and
\lean{Module.Finite.toModuleEnd_moduleEnd_surjective}.
The local lemmas
\lean{exists_trace_action_ne_of_not_linearEquiv} and
\lean{linearMap_trace_surjective}, in the same namespace, make the
separating operator explicit. The theorem itself applies to a group over
an algebraically closed field of any characteristic and does not require
the group to be finite. Its use here is in characteristic zero.
\end{implementation}

For a finite group, equality of irreducible characters also implies
equivalence over an arbitrary field $K$ of characteristic zero. The
character averaging formula (see, for example,
\cite[Section~12.1, proof of Proposition~32]{Serre1977}) gives
\[
 \dim_K\operatorname{Hom}_{KG}(V,W)
   =\frac{1}{|G|}\sum_{g\in G}\chi_V(g^{-1})\chi_W(g),
\]
where the integer on the left is viewed in $K$. If $\chi_V=\chi_W$,
the right hand side is also $\dim_K\operatorname{End}_{KG}(V)$.
This dimension is positive because $V\ne0$ and its identity map is
nonzero. Thus there is a nonzero homomorphism $V\to W$, which is an
isomorphism by irreducibility. The proof does not require that the
endomorphism ring consist of scalars.

\begin{implementation}
The theorem
\lean{ModularRep.PaperProofs.TypeBOrdinaryTraceSeparationSplitting.nonempty_equiv_of_character_eq}
proves this implication using Mathlib's character averaging formula. It applies
to every finite group over every field of characteristic zero and also
identifies simple module classes with equal characters. In particular,
it applies when $K$ is the fraction field of a discrete valuation ring.
\end{implementation}

\section{Pullback and quotient descent}
\label{lib:ordinary:pullback}

For a homomorphism $f:H\to G$, the pullback $f^*\rho$ acts on the same
vector space by $(f^*\rho)(h)=\rho(f(h))$. Its character is
$\chi_\rho\circ f$. An intertwiner between two representations of $G$
also intertwines their pullbacks.

\begin{proposition}[Inflation and descent of representations]
\label{lib:ordinary:surjective-pullback}
If $f:H\twoheadrightarrow G$ is surjective, then $f^*\rho$ is irreducible
if and only if $\rho$ is irreducible. If $\pi:G\twoheadrightarrow H$ and
$\ker \pi\leq\ker\rho$, there is a unique representation $\bar\rho$ of $H$
on $V$ such that $\pi^*\bar\rho=\rho$. It is irreducible whenever $\rho$ is.
\end{proposition}
\begin{proof}
A subspace invariant under $\rho$ is invariant under its pullback. Conversely,
surjectivity lets every $g\in G$ be written as $f(h)$, so an invariant
subspace for the pullback is invariant under $\rho$. The lattices of
invariant subspaces are identical.

For descent, define $\bar\rho(\pi(g))=\rho(g)$. If $\pi(g)=\pi(g')$, then
$g'^{-1}g\in\ker \pi\leq\ker\rho$, which makes the value independent of the
lift. This defines a representation and proves the pullback identity.
Surjectivity proves uniqueness, and the first assertion proves
irreducibility.
\end{proof}

\begin{implementation}
\module{ModularRep.Twist} constructs the order isomorphism
\lean{Representation.subrepresentationPullbackOrderIso}, from which
\lean{Representation.isIrreducible_pullback_iff} follows. These results
hold over any field, and the order isomorphism is available more generally
over a semiring. In
\module{ModularRep.RepresentationSurjectiveDescent},
\lean{descend} factors through $G/\ker \pi$ and uses Mathlib's quotient
isomorphism $G/\ker \pi\simeq H$. The declarations
\lean{descend_apply}, \lean{descend_pullback} and
\lean{descend_irreducible} give its defining identities and irreducibility.
Uniqueness on the fixed vector space is the direct surjectivity argument
in the proof above.
\end{implementation}

To descend a character, it remains to show that the kernel of the quotient
map acts trivially in an affording representation. Over $\mathbb C$, the
standard identity
$\ker\rho=\{g\in G:\chi_\rho(g)=\chi_\rho(1)\}$ gives this directly
(see, for example, \cite[Lemma~2.19]{Isaacs1976}). The averaging argument below proves
the needed implication over any field of characteristic zero, without
assuming irreducibility.

\begin{lemma}[Kernel containment by averaging]
\label{lib:ordinary:averaging}
Let $S$ be a finite subgroup of $G$. If
$\chi_\rho(s)=\chi_\rho(1)$ for every $s\in S$, then $S\leq\ker\rho$.
\end{lemma}
\begin{proof}
The averaging operator
\[
 P_S=\frac1{|S|}\sum_{s\in S}\rho(s)
\]
is a projection onto the invariant subspace $V^S$. Its trace is both
$\dim_K V^S$ and, by the hypothesis, $\dim_K V$. These integers are equal
because their images in a field of characteristic zero are equal. Hence
$V^S=V$, and every element of $S$ acts trivially.
\end{proof}

\begin{implementation}
\lean{Representation.subgroup_le_ker_of_character_eq_one}, in
\module{ModularRep.OrdinaryCharacterLiesOverKernel}, proves this using Mathlib's
\lean{Representation.card_inv_mul_sum_char_eq_finrank} for the restricted
representation. Only $S$ must be finite. No irreducibility, normality or
centrality hypothesis occurs in this lemma.
\end{implementation}

\begin{theorem}[Descent of an ordinary irreducible character]
\label{lib:ordinary:character-descent}
Let $\pi:G\twoheadrightarrow H$ be a surjective homomorphism and let
$\chi\in\Irr_K(G)$. Suppose that $\chi(x)=\chi(1)$ for all
$x\in\ker \pi$. Then there is a unique $\bar\chi\in\Irr_K(H)$ with
\[
 \chi(g)=\bar\chi(\pi(g))\qquad(g\in G).
\]
\end{theorem}
\begin{proof}
Choose an irreducible realisation of $\chi$. Lemma~\ref{lib:ordinary:averaging}
places $\ker \pi$ in its kernel. Descend this representation by
Proposition~\ref{lib:ordinary:surjective-pullback} and take its trace.
The resulting representation remains irreducible. Any two functions with
the required pullback agree at every element of $H$, since $\pi$ is surjective.
\end{proof}

\begin{implementation}
The descended character, its factorisation and its uniqueness are given by
\lean{ModularRep.OrdinaryIrreducibleCharacterSurjectiveDescent.descendCharacter},
\lean{descendCharacter_factorisation} and \lean{descendCharacter_unique}.
The construction applies averaging to an irreducible realisation and
then descends that representation.
\end{implementation}

For a positive integer $a$, write $a_\ell$ for the largest power of the
prime $\ell$ dividing $a$. An ordinary irreducible character $\chi$ has
\emph{defect zero at $\ell$} if $\chi(1)_\ell=|G|_\ell$, where
$\chi(1)$ denotes its positive integral degree.
Over a splitting field this is the usual notion of defect zero. Over an
arbitrary field of characteristic zero, we use the same numerical
condition on the dimension of an irreducible representation in the
following statement.

\begin{corollary}[Preservation of defect zero]
\label{lib:ordinary:defect-descent}
Let $\ell$ be a prime. Under the hypotheses of
Theorem~\ref{lib:ordinary:character-descent}, assume that $\ell\nmid|\ker \pi|$.
If $\chi$ has defect zero at $\ell$, then $\bar\chi$ has defect zero.
\end{corollary}
\begin{proof}
An affording representation $V$ has dimension $\chi(1)$. Descent retains
this vector space and hence the degree.
The equality $|G|=|H|\,|\ker \pi|$ gives $|G|_\ell=|H|_\ell$, so the same
dimension satisfies the equality defining defect zero for $H$.
\end{proof}

\begin{implementation}
\lean{ModularRep.IsDefectZeroOrdinaryCharacter} expresses the dimension
equality for an irreducible representation in Mathlib's \lean{FDRep}
affording the character.
\lean{descendCharacter_defectZero} in the descent namespace applies
averaging to the affording representation and descends it without changing
its dimension.
\end{implementation}

\begin{example}[Projection from a direct product]
Let $G=H\times C_r$, with $\ell\nmid r$, and let $\pi$ be the first
projection. To apply the descent theorem to an ordinary irreducible
character $\chi$, it is enough to verify
$\chi(1,c)=\chi(1,1)$ for every $c\in C_r$. The theorem then constructs
$\bar\chi\in\Irr_K(H)$ and proves
$\chi(h,c)=\bar\chi(h)$ for all $(h,c)$. If the original character has
defect zero at $\ell$, the corollary preserves that property.
\end{example}

\section{Linear characters and automorphisms}
\label{lib:ordinary:actions}

Write $L_K(G)=\operatorname{Hom}(G,K^\times)$ for the group of linear
characters. For $\lambda\in L_K(G)$ define
$(\rho\otimes\lambda)(g)=\lambda(g)\rho(g)$ on $V$. The usual tensor
product with a one dimensional representation is realised on $V$ by
this formula. For an automorphism $\alpha$ of $G$, set
$\rho^\alpha=\rho\circ\alpha$, $\lambda^\alpha=\lambda\circ\alpha$
and $\chi^\alpha=\chi\circ\alpha$.

\begin{proposition}[Tensoring with a linear character]
\label{lib:ordinary:linear-twist}
Tensoring with a linear character preserves and reflects irreducibility,
and
\[
 \chi_{\rho\otimes\lambda}(g)=\lambda(g)\chi_\rho(g).
\]
For $\alpha\in\Aut(G)$, one also has
\[
 (\rho\otimes\lambda)^\alpha
   =\rho^\alpha\otimes\lambda^\alpha,
 \qquad
 (\chi\lambda)^\alpha=\chi^\alpha\lambda^\alpha.
\]
\end{proposition}
\begin{proof}
Multiplication by the nonzero scalar $\lambda(g)$ does not change which
subspaces are invariant under the action of $g$. Multiplication by
$\lambda(g)^{-1}$ gives the converse. The trace formula follows from
linearity of trace. The compatibility with $\alpha$ follows by evaluating
both sides at $g$.
\end{proof}

\begin{implementation}
\module{ModularRep.LinearCharacterTensorAction} defines
\lean{Representation.linearCharacterTwist} and proves
\lean{Representation.isIrreducible_linearCharacterTwist_iff},
\lean{Representation.character_linearCharacterTwist} and
\lean{Representation.linearCharacterTwist_twist}.
The induced operation on ordinary irreducible characters is
\lean{ModularRep.OrdinaryIrreducibleCharacter.linearTwist}.
\end{implementation}

The convention for automorphisms is a right action:
$(\chi^\alpha)^\beta=\chi^{\alpha\beta}$, with
$(\alpha\beta)(g)=\alpha(\beta(g))$. The combined tensor and automorphism
action below is instead written as a left action, with the indicated inverses.

More precisely, let $A$ act on $G$ through $a\mapsto\alpha_a$. Its left
action on $L_K(G)$ is $a\cdot\lambda=\lambda\circ\alpha_{a^{-1}}$.
The resulting semidirect product acts on $\Irr_K(G)$ by
\[
 (\lambda,a)\cdot\chi
   =\lambda^{-1}\,\chi^{\alpha_{a^{-1}}}.
\]
The compatibility in Proposition~\ref{lib:ordinary:linear-twist} proves
the action law. The inverses affect formulas for the action, although they
do not change the orbits or stabilisers of either factor acting alone.

\begin{implementation}
The declarations
\lean{ModularRep.OrdinaryIrreducibleCharacter.linearCharacterFieldAction},
\lean{tensorField_semidirectCompatible},
\lean{tensorFieldSemidirectAction} and
\lean{tensorFieldSemidirectAction_apply} specify this convention.
They use the semidirect product action in
\module{Formalisation.SemidirectStabilizer}. The right action by
automorphisms is represented as a left action of the opposite group.
These declarations apply to any group $A$ acting on $G$.
\end{implementation}

\begin{example}[A stabiliser calculation]
Suppose $A$ and a subgroup $L\leq L_K(G)$ are given, and $A$ preserves $L$.
Under the combined action, $(\lambda,a)$ fixes $\chi$ if and only if
$\chi(\alpha_{a^{-1}}(g))=\lambda(g)\chi(g)$ for every $g\in G$.
Pullback followed by tensoring with a linear character gives an irreducible
representation affording the transformed character. To restrict this
action to the characters in a block, one must also prove that the block
is preserved.
\end{example}

\section{Extension across a cyclic quotient}
\label{lib:ordinary:cyclic-extension}

Let $N\trianglelefteq H$ with $H$ finite. An extension of a representation
$\rho$ of $N$ is a representation $\widehat\rho$ of $H$ on the same vector
space such that
$\widehat\rho|_N\simeq\rho$. An extension makes $\rho$ invariant
under conjugation by $H$: the action of $h^{-1}$ intertwines the
conjugation twist of $\widehat\rho|_N$ with $\widehat\rho|_N$.
Also, an extension of an irreducible representation is irreducible, since
every $H$-invariant subspace is $N$-invariant.

\begin{theorem}[Cyclic extension]
\label{lib:ordinary:cyclic-principle}
Assume that $K$ is algebraically closed. Let $H$ be finite, let
$N\trianglelefteq H$, and let $\rho$ be a finite
dimensional irreducible $K$-representation of $N$. If $H/N$ is cyclic
and $\rho\circ c_h\simeq\rho$ for every $h\in H$, where
$c_h(n)=hnh^{-1}$, then $\rho$ extends to $H$.
\end{theorem}
\begin{proof}
Choose $h\in H$ whose coset generates $H/N$, and set $r=|H/N|$.
Conjugation invariance gives an invertible operator $T$ such that
$T\rho(n)T^{-1}=\rho(hnh^{-1})$ for every $n\in N$.
Consequently $\rho(h^r)^{-1}T^r$ commutes with $\rho(N)$ and is
a nonzero scalar $c$ by Schur's lemma. Choose $u\in K$ with
$u^r=c^{-1}$. Replacing $T$ by $uT$ preserves the conjugation
identity and gives $T^r=\rho(h^r)$.
These identities ensure that $\widehat\rho(nh^i)=\rho(n)T^i$,
for $n\in N$ and $0\leq i<r$, defines a representation of $H$
extending $\rho$.
\end{proof}

For the complex character formulation, see
\cite[Corollary~11.22]{Isaacs1976}. The proof above uses
algebraic closure to apply Schur's lemma and take the required root.
The modular analogue is discussed in
Section~\ref{lib:brauer:actions-extensions}.

\begin{proposition}[Extension of an invariant ordinary character]
\label{lib:ordinary:cyclic-bridge}
Assume that $K$ is algebraically closed. Let $H$ be finite, let
$N\trianglelefteq H$, let $H/N$ be cyclic, and let $\chi\in\Irr_K(N)$. If
\[
 \chi(hnh^{-1})=\chi(n)\qquad(h\in H,\ n\in N),
\]
then there exists $\widehat\chi\in\Irr_K(H)$ such that
$\widehat\chi|_N=\chi$.
\end{proposition}
\begin{proof}
Choose an irreducible realisation $\rho$ of $\chi$. Each conjugation twist
of $\rho$ has character $\chi$, so
Proposition~\ref{lib:ordinary:separation} supplies an equivalence with
$\rho$. The cyclic extension theorem gives $\widehat\rho$ and its
restriction equivalence. The extension is irreducible, and taking traces
of the intertwining identity proves the required restriction equality.
\end{proof}

\begin{implementation}
\module{ModularRep.CyclicExtension} defines
\lean{Representation.ConjugationInvariant} and
\lean{Representation.Extension}.
The Lean construction assumes
\lean{Representation.CyclicExtensionPrinciple}, which states the
representation extension property over the chosen coefficient field.
This principle is an explicit hypothesis of the Lean construction,
including when $K$ is algebraically closed. Theorem~\ref{lib:ordinary:cyclic-principle}
proves the mathematical assertion under algebraic closure.
Over an arbitrary field of characteristic zero the assertion can fail. In
\module{ModularRep.OrdinaryCharacterCyclicExtensionBridge},
\lean{ModularRep.OrdinaryIrreducibleCharacter.realisation_conjugationInvariant_of_fixed}
proves invariance of an irreducible realisation, and
\lean{ModularRep.OrdinaryIrreducibleCharacter.exists_extension_of_fixed_cyclic_quotient}
constructs the extending character under the stated extension hypothesis.
An \lean{ExtensionWitness} consists of
an irreducible character of $H$ together with its restriction identity
on $N$.
\end{implementation}

\begin{example}[Extension to a semidirect product]
Let $D$ and $A$ be finite, with $A$ cyclic. For $H=D\rtimes A$, the normal
copy of $D$ has quotient $H/D\simeq A$. A conjugation invariant irreducible
representation of $D$ therefore extends to $H$ whenever the cyclic
extension property holds over the coefficient field.
\lean{Representation.exists_extension_to_semidirect_of_cyclic_outer}
constructs this quotient isomorphism and applies the principle.
\end{example}
